\section{Extremal selection and propagation}\label{sec:extremal}

We prove Theorem~\ref{pkt:main}. We use the packet scales, lattice norms, and
propagation exponent introduced in Section~\ref{sec:packets}. In particular,
the amplification of a diagram with initial array \(d\) and terminal array
\(c\) is
\[
                 \frac{\mathcal A(c)^3}{G(d)^3}.
\]
Every mask below is a coordinate projection on a layer of the diagram.
Between successive masks we retain the entire packet transition matrix.
Graphs will record its significant interactions, but will not replace the
matrix by an edgewise truncation.

We argue by contradiction and assume \(\gamma>10\kappa\).
We use two successive extrema. The first determines the infimum of
the time-support exponents that allow amplification arbitrarily close
to \(\gamma\). The second minimizes, up to a vanishing error, the
asymptotic exponent of the path count per original root among diagrams
with nearly this time support. This makes the path count stable under
the later masks that retain near-extremal amplification, and supplies
the counting law needed for geometric propagation.

\subsection{The least time dimension of a near-extremal family}

For \(0\leq \sigma\leq1\), define \(\gamma_\sigma\) by the same supremum
of limiting amplification exponents as \(\gamma\), with the following
additional restriction. There is a fixed \(\eta>0\) within the family
such that each active initial index touches terminal masked indices in
at most \(m^{2\sigma+o(1)}\) time bins, where a touch means
\[
 |\nu_b-\nu_v|\leq m^\eta\theta,
 \qquad
 |z_b-z_v-t_{I_b}U_{\nu_v}|\leq m^\eta\delta .
\]
Here the initial time is zero. The supremum still ranges over all fixed
positive \(\eta\) and all finite polynomial parameter bounds.

\begin{lemma}[Minimal time dimension]\label{ext:time-dimension}
One has
\begin{equation}\label{eq:source-34}
 \gamma_\sigma\leq 2\sigma,
 \qquad
 s:=\inf\{\sigma\in[0,1]:\gamma_\sigma=\gamma\}>0 .
\end{equation}
For every fixed \(\sigma<s\), the difference
\(\gamma-\gamma_\sigma\) is positive. For every \(\sigma>s\) in
\([0,1]\), one has \(\gamma_\sigma=\gamma\).
Near-extremal families defining any of these exponents may be chosen
with all nonzero initial coordinate amplitudes between \(1\) and \(2\).
\end{lemma}

\begin{proof}
Fix a family satisfying the additional support restriction.
For a terminal bin \(I\), let \(\mathcal R_I\) be the initial indices
touching at least one retained terminal index in \(I\). Every matrix
entry from an initial index outside \(\mathcal R_I\) to a retained
index in \(I\) is a locality tail. Lemma~\ref{pkt:locality} and the
polynomial coordinate count show that replacing the input by
\(d|_{\mathcal R_I}\) changes the output in \(I\) by a negligible
error. Apply the squared-sum contraction of Lemma~\ref{pkt:crude} to
the full transition matrix on this restricted input. Summing over
\(I\) counts each initial index at most \(m^{2\sigma+o(1)}\) times.
Therefore
\[
 r w^2\sum_b |c_b|^2
 \leq m^{-2(1-\sigma)+o(1)}\delta^2\sum_v|d_v|^2.
\]
The same lemma gives \(K(c_I)\leq m^{4+o(1)}K(d)\) in every bin.
Since \(A(c_I)\leq G(c_I)\), multiplication by the square root of this
capacity bound proves amplification at most \(m^{2\sigma+o(1)}\).

The classes of admissible families increase with \(\sigma\).
Consequently \(\gamma_\sigma\) is nondecreasing and at most \(\gamma\);
also \(\gamma_1=\gamma\), since there are only \(m^2\) terminal time
bins. The set on which equality holds is therefore an upper interval,
possibly without its lower endpoint. Its infimum satisfies
\(s\geq\gamma/2>0\), and the other assertions about
\(\gamma_\sigma\) follow. Equality at \(\sigma=s\) is not needed.

To obtain comparable initial amplitudes, first normalize the largest
one. The polynomial coordinate counts and the polynomial comparison
between coordinate norms and the packet norms allow amplitudes below
a sufficiently small fixed negative power of \(m\) to be discarded.
The remaining amplitudes lie in \(O(\log m)\) dyadic classes.
Writing the input as their sum and using linearity followed by the
triangle inequality for the output norm, one class retains at least
an inverse-logarithmic fraction of the output norm. Restriction to
that class can only decrease the initial raw gauge. Its cubed
amplification therefore loses a subpower factor. Rescaling its
amplitudes gives the interval \([1,2]\). The support restriction is
hereditary under this operation.
\end{proof}

An upper bound with exponent \(\gamma_\sigma\) also holds with the
initial atomic norm in place of the raw gauge: decompose the input
into atoms on its given support and apply the triangle inequality.
In particular, Lemma~\ref{pkt:composition} permits this improved bound
at one transition of a masked diagram whenever the two supports of
that transition satisfy the time restriction. We will repeatedly use
the following consequence:
\begin{equation}\label{ext:gap}
 \begin{gathered}
 \text{a fixed improvement below the time exponent \(s\) at one transition}\\
 \text{gives a fixed improvement below \(\gamma\) for the diagram.}
 \end{gathered}
\end{equation}
The improvement in \eqref{ext:gap} is measured in \(m\)-power units
when the transition ratio is a fixed positive power of \(m\).

\subsection{A second extremum: the number of paths}

Fix \(J\geq1\), and let \(m_j\), \(0\leq j\leq J\), be dyadic
roundings of \(m^{j/J}\), with \(m_0=1\) and \(m_J=m\). Put
\[
 \ell_j=\frac{m_{j+1}}{m_j},\qquad
 \theta_j=m_j\delta,\qquad r_j=m_j^{-2},\qquad w_j=r_j\theta_j .
\]
Thus \(\log_m\ell_j\to1/J\). Insert the intermediate frame
identities and truncate to polynomial grids, as permitted by
Lemma~\ref{pkt:locality}. A diagram consists of these exact transitions
and arbitrary vertex masks, including a possible initial mask.
Its input array \(d\) is normalized as in Lemma~\ref{ext:time-dimension},
on a root set \(\mathcal R\); amplification is always measured against
the entire array on this set, even if an additional root mask is used.

For each fixed \(J\), an index \(i\) will label a sequence of families.
Within family \(i\), the polynomial bounds of
Definition~\ref{pkt:families} and all positive radius powers are fixed
before \(m\to\infty\); these bounds may depend on \(i\).
We write \(o_i(1)\) in an exponent for an error whose limiting upper
absolute value as \(m\to\infty\) tends to zero as \(i\to\infty\).
This includes subpower losses within a fixed family. Thus \(J\)
remains fixed while we first take the inner limit in \(m\) and then
the outer limit in \(i\). The construction is valid for every fixed
\(J\); we will later choose one sufficiently large for the geometric
time-profile hypothesis.

For a fixed positive radius power \(\eta\), join an active index \(v\)
at level \(j\) to an active index \(b\) at level \(j+1\) when
\begin{equation}\label{eq:source-35}
 I_b\subset I_v,\qquad
 |\nu_b-\nu_v|\leq m^\eta\theta_{j+1},\qquad
 |z_b-z_v-(t_{I_b}-t_{I_v})U_{\nu_v}|\leq m^\eta w_j .
\end{equation}
A path is a sequence of such edges from an active root to a terminal
leaf.

\paragraph{Locality for error diagrams.}\label{ext:diagram-locality}
We will repeatedly discard a vertex class whose contribution has a
nonadjacent edge in every complete matrix-product summand.
For every prescribed \(A>0\), this contribution is
\(O(m^{-A}\max_{v\in\mathcal R}|d_v|)\), both coordinatewise and in
the packet norms. Indeed, Lemma~\ref{pkt:locality} makes the
nonadjacent coefficient smaller than any prescribed negative power;
the remaining coefficients and the number of summands have polynomial
bounds, and there are only \(J\) layers. Choose the locality decay
order after the graph radius and these fixed family bounds.
This applies, in particular, to roots with no graph path to the tested
output and to intermediate vertices with no suffix to that output.
It bounds the error of the vertex masking; the retained diagram
continues to use the full transition matrices.

\begin{lemma}[Endpoint localization]\label{ext:endpoints}
Along a path from level \(a\) to level \(b>a\), the endpoint angle
difference is \(O_J(m^\eta\theta_b)\), and
\[
 |z_b-z_a-(t_{I_b}-t_{I_a})U_{\nu_a}|
       \leq O_J(m^\eta w_a).
\]
For a fixed root and a fixed terminal leaf, the number of paths is
at most \(m^{O(J\eta)+o(1)}\).
\end{lemma}

\begin{proof}
The angular differences are bounded by a geometric sum of the
successive child widths. For the position bound, write
\(t_j=t_{I_j}\), \(U_j=U_{\nu_j}\), and
\(e_j=z_{j+1}-z_j-(t_{j+1}-t_j)U_j\). Telescoping gives
\[
 z_b-z_a-(t_b-t_a)U_a
 =\sum_{j=a}^{b-1}e_j+
   \sum_{j=a+1}^{b-1}(t_b-t_j)(U_j-U_{j-1}).
\]
The edge errors satisfy \(|e_j|\leq m^\eta w_j\).
For a switch from level \(j-1\) to level \(j\), Lipschitz
continuity of the velocity map gives
\(|U_j-U_{j-1}|\lesssim m^\eta\theta_j\), while nesting gives
\(0\leq t_b-t_j\leq r_j\). Its contribution is therefore
\(O(m^\eta\theta_jr_j)=O(m^\eta w_j)\).
Summing and using \(w_j\leq w_a\) proves the positional endpoint bound.

For fixed endpoints, the intermediate time interval at each level is
the unique ancestor of the terminal interval. At level \(j\), the
angular grid has spacing \(\theta_j\), while the angle lies within
\(O_J(m^\eta\theta_j)\) of the root angle. There are
\(m^{O(\eta)+o(1)}\) choices. Given this angle, the endpoint bound
applied from level \(j\) to the terminal leaf locates \(z_j\) within
\(O_J(m^\eta w_j)\) of a prescribed point. On the position grid of
spacing comparable to \(w_j\), this again gives
\(m^{O(\eta)+o(1)}\) choices. Multiplication over the fixed number
of intermediate levels proves the claim.
\end{proof}

At accuracy \(\epsilon>0\), call a family of diagrams admissible if it
has a fixed radius power \(0<\eta<\epsilon\), amplification at least
\(m^{\gamma-\epsilon}\), and at most \(m^{2(s+\epsilon)}\) terminal
time bins accessible from every root by paths at that radius.
For such a family consider the limiting lower exponent of
\begin{equation}\label{eq:source-36}
                 \frac{\#\{\text{paths}\}}{\#\mathcal R}.
\end{equation}

\Needspace{4\baselineskip}
\begin{lemma}[Minimal path count]\label{ext:path-minimum}
Admissible families exist at arbitrarily small accuracies. The infima
of the limiting lower exponents in \eqref{eq:source-36} have a finite
limit as \(\epsilon\downarrow0\). There are families approaching this
limit, with accuracies tending to zero, for which the following holds.
Call their graphs the high graphs. Choose a smaller fixed positive
radius power in each family, tending to zero with its accuracy, and
call the corresponding graphs low graphs. After any fixed number of
further vertex maskings that retain amplification
\(m^{\gamma-o_i(1)}\), one has
\begin{equation}\label{eq:source-37}
 \#\{\text{remaining low paths}\}
       \geq m^{-o_i(1)}\#\{\text{original high paths}\}.
\end{equation}
\end{lemma}

\begin{proof}
If \(s<1\), choose a support exponent strictly between \(s\) and
\(s+\epsilon\); its propagation exponent is \(\gamma\) by
Lemma~\ref{ext:time-dimension}. If \(s=1\), use the unrestricted
families. In either case choose a near-extremal single-step family
with the desired time sparsity at some fixed enlarged direct radius,
and make its input amplitudes comparable. Insert the intermediate
identities. By Lemma~\ref{ext:endpoints}, a sufficiently smaller fixed
radius power at the intermediate levels ensures that every path
endpoint is a direct touch at the originally chosen radius.
Truncation errors and dyadic losses can be made smaller than the
unused accuracy. This gives the required admissible families.

At each step, a parent has at most
\(O(m^{4\eta}\ell_j^2)\) children in a given time bin: the angular
grid contributes \(O(m^{2\eta})\) choices and the position grid
contributes \(O(m^{2\eta}\ell_j^2)\). There are \(\ell_j^2\) child
time bins. Hence the total path count per root is at most
\(m^{4+4J\eta+o(1)}\). For a lower bound, let
\(\mathcal R'\) be the roots having at least one path. All other
roots contribute a negligible error by the
\hyperref[ext:diagram-locality]{locality observation above}.
Since \(1\leq|d_v|\leq2\),
\[
 G(d|_{\mathcal R'})^3
 \leq C\,\frac{\#\mathcal R'}{\#\mathcal R}\,G(d)^3.
\]
Here the squared mass drops in the displayed proportion, and its
capacity does not increase. Composition with exponent
\(\gamma+o(1)\), together with the lower amplification, gives
\(\#\mathcal R'/\#\mathcal R\geq m^{-\epsilon-o(1)}\).
There is at least one path per root of \(\mathcal R'\), proving
a lower bound for the exponent in \eqref{eq:source-36}.
The admissible classes are nested as \(\epsilon\) decreases, so
their infima have a finite monotone limit; denote it by \(\Lambda_J\).

Choose near-minimizers with accuracies \(\epsilon_i\downarrow0\),
and pass to subsequences on which their path-count exponents
approach \(\Lambda_J\) also from above. Let their fixed high radius
powers be \(\eta_i\), and choose low powers
\(0<\eta_i^{\mathrm{lo}}\ll\eta_i\).
Suppose further masks have amplification \(m^{\gamma-\rho_i+o(1)}\)
with \(\rho_i\to0\). The resulting low graph inherits the original
upper bound on accessible terminal bins. It is therefore an
admissible competitor with accuracy tending to zero, for example
an accuracy slightly larger than
\(\max(\epsilon_i,\rho_i,\eta_i^{\mathrm{lo}})\).
Its limiting lower exponent is at least \(\Lambda_J-o_i(1)\).
The original high graph has exponent at most
\(\Lambda_J+o_i(1)\) on the chosen subsequence. This proves
\eqref{eq:source-37}. All adaptive masks still form families with
fixed polynomial bounds at each \(i\), so this comparison applies
to each of a fixed finite number of such selections.
\end{proof}

Figure~\ref{fig:packet-proof-map} shows how this path-retention property
enters the rest of the proof.
\begin{figure}[tb]
\centering
\begin{tikzpicture}[
  >=Stealth,
  proofbox/.style={draw=linkblue!70,fill=linkblue!4,rounded corners=2pt,
    align=center,text width=4.6cm,minimum height=1.05cm,
    inner sep=5pt,font=\small},
  proofarrow/.style={->,thick,draw=linkblue!80}]
  \node[proofbox] (time) at (-3.15,0)
    {First extremum: \(s\)\\\(\gamma_\sigma<\gamma\) for \(\sigma<s\)};
  \node[proofbox] (paths) at (3.15,0)
    {Second extremum: paths per root\\retain paths under later masks};
  \node[proofbox] (law) at (-3.15,-1.75)
    {Control conditional time-bin\\counts and probabilities};
  \node[proofbox] (leaves) at (3.15,-1.75)
    {Select terminal coefficients\\control incoming path counts};
  \node[proofbox] (capacity) at (-3.15,-3.5)
    {Geometric propagation\\bound output capacity};
  \node[proofbox] (mass) at (3.15,-3.5)
    {Exact maps and refined decoupling\\second and fourth moments};
  \node[proofbox,text width=5.2cm] (end) at (0,-5.3)
    {Capacity--mass cancellation\\\(\gamma\le10\kappa\)};
  \draw[proofarrow] (time)--(paths);
  \draw[proofarrow] (paths)--(leaves);
  \draw[proofarrow] (leaves)--(law);
  \draw[proofarrow] (law)--(capacity);
  \draw[proofarrow] (leaves)--(mass);
  \draw[proofarrow] (capacity.south)--(end.north west);
  \draw[proofarrow] (mass.south)--(end.north east);
\end{tikzpicture}
\caption{The two extrema in the packet proof. Here \(s\) is the infimum
of the terminal time-support exponents permitting amplification approaching
\(\gamma\). Graph paths run from initial packet indices (roots) to
terminal indices (leaves). Minimizing the limiting exponent of the
path count per original root
forces later coordinate selections to retain enough paths whenever
they preserve nearly maximal amplification. Uniform counting on a
selected path set supplies the time law needed for geometric
propagation. The second- and fourth-moment estimates concern the
exact masked coefficient maps.}
\label{fig:packet-proof-map}
\end{figure}


Fix these high and low radii. Between them choose a fixed intermediate
radius power \(\eta_i^{\mathrm{mid}}\) in each family, so that
\(\eta_i^{\mathrm{lo}}<\eta_i^{\mathrm{mid}}<\eta_i\).
The strict gaps absorb bounded constants in neighborhood containments
once \(m\) is sufficiently large. Figure~\ref{fig:adjacency-radii}
illustrates the inclusions used below.

\begin{figure}[tb]
\centering
\begin{tikzpicture}[x=1cm,y=1cm,>=Stealth,font=\small]
  \draw[gray!65,->] (-.4,0)--(10.7,0) node[right] {time};
  \draw[gray!65,->] (0,-1.55)--(0,1.7) node[above] {position};
  \fill[linkblue!4] (0,-1.2)--(10,-.2)--(10,2.2)--(0,1.2)--cycle;
  \fill[linkblue!17] (0,-.4)--(10,.6)--(10,1.4)--(0,.4)--cycle;
  \draw[linkblue,thick] (0,0)--(10,1);
  \draw[linkblue!50,dashed] (0,-1.2)--(10,-.2);
  \draw[linkblue!50,dashed] (0,1.2)--(10,2.2);
  \draw[linkblue!80] (0,-.4)--(10,.6);
  \draw[linkblue!80] (0,.4)--(10,1.4);
  \draw[densely dotted] (5,-1.15)--(5,1.8);
  \fill (5,.66) circle (2pt)
    node[right=4pt,above=1pt] {\(b\)};
  \fill (5,.30) circle (2pt)
    node[left=4pt,below=1pt] {\(b'\)};
  \node[anchor=west,fill=white,inner sep=2pt] at (7.5,1.75)
    {parent high};
  \node[anchor=west,fill=white,inner sep=2pt] at (7.8,1.02)
    {parent low};
  \node[anchor=north] at (5,-1.28) {one child time bin};
  \node[anchor=south west] at (.6,.15) {parent centerline};

  \node at (5,-2.10) {Position cross-section at the marked bin};
  \begin{scope}[shift={(5,-3.55)},x=1.3cm]
    \filldraw[fill=linkblue!4,draw=linkblue!50,dashed]
      (-3,-.42) rectangle (3,.42);
    \filldraw[fill=linkblue!10,draw=linkblue!65]
      (-1.6,-.30) rectangle (2.4,.30);
    \filldraw[fill=linkblue!17,draw=linkblue!80]
      (-1,-.18) rectangle (1,.18);
    \draw[gray!65,densely dotted] (0,-.90)--(0,.50);
    \draw[linkblue!65,densely dotted] (.4,-.30)--(.4,.30);
    \fill (.4,0) circle (2pt);
    \fill (-.5,0) circle (2pt);
    \node[anchor=north] at (.4,-.48) {\(b\)};
    \node[anchor=north] at (-.5,-.48) {\(b'\)};
    \node[anchor=north,font=\scriptsize,text=gray!80] at (0,-.96)
      {parent center};
    \node at (-2.25,.80) {parent high};
    \draw[linkblue!50] (-2.25,.59)--(-2.55,.42);
    \node at (1.10,.80) {intermediate about \(b\)};
    \draw[linkblue!65] (1.10,.59)--(1.65,.30);
    \node at (-2.20,-.72) {parent low};
    \draw[linkblue!80] (-1.60,-.62)--(-1,-.18);
  \end{scope}
\end{tikzpicture}
\caption{Schematic localization neighborhoods in one position coordinate.
At the marked child time bin, the intermediate neighborhood centered at the
selected low child \(b\) contains every low child of the parent in that bin
and is contained in the parent's high neighborhood. The angular bounds in
the actual neighborhoods are not depicted. These inclusions are used for
counting paths; the analytic evolution uses the exact coefficient maps with
vertex masks.}
\label{fig:adjacency-radii}
\end{figure}


\subsection{Regularizing the graph by vertex masks}

We first record how dyadic selections affect the exact diagram.
Partitioning the active coordinates at any one layer into
\(O((\log m)^C)\) classes writes the exact output as a sum of the
outputs with the corresponding vertex masks inserted. The triangle
inequality for \(\mathcal A\) selects a class losing only a subpower
factor in amplification. This remains true through finitely many
layers. The integer path statistics below have polynomial size, so
their positive values require only logarithmically many dyadic
classes. Vertices with no high suffix have negligible output
contribution by the
\hyperref[ext:diagram-locality]{locality observation} and can be discarded.

Let \(D_j(v)\) be the number of suffix paths from \(v\) in the
\emph{original high graph}. For a vertex \(b\) at level \(j+1\),
let \(\mathcal N_j(b)\) consist of the original high vertices \(b'\)
at level \(j+1\) satisfying
\[
 I_{b'}=I_b,\qquad
 |\nu_{b'}-\nu_b|\leq m^{\eta_i^{\mathrm{mid}}}\theta_{j+1},
 \qquad
 |z_{b'}-z_b|\leq m^{\eta_i^{\mathrm{mid}}}w_j.
\]
Define
\[
 S_j(b)=\sum_{b'\in\mathcal N_j(b)}D_{j+1}(b').
\]
If \(b\) is a low child of \(v\), the triangle inequality shows that
\(\mathcal N_j(b)\) contains every low child of \(v\) in that time
bin and is contained in the high children of \(v\). Indeed, for fixed
\(i\) and sufficiently large \(m\),
\[
 2m^{\eta_i^{\mathrm{lo}}}\leq m^{\eta_i^{\mathrm{mid}}},
 \qquad
 m^{\eta_i^{\mathrm{mid}}}+m^{\eta_i^{\mathrm{lo}}}\leq m^{\eta_i}.
\]
The predicted child position in \eqref{eq:source-35} is the same
for all children in that bin.

\Needspace{18\baselineskip}
\begin{lemma}[Time branching and pruning]\label{ext:branching}
One can retain amplification \(m^{\gamma-o_i(1)}\) by vertex masks such that
\[
                         D_j(v)\asymp d_j,\qquad
                         S_j(b)\asymp s_j'
\]
on the selected coordinates, for \(0\leq j\leq J\) in the
first comparison and \(0\leq j<J\) in the second, with
\begin{equation}\label{eq:source-38}
                         \frac{d_j}{s_j'}
                    \geq \ell_j^{2s}m^{-o_i(1)}.
\end{equation}
One can then prune backwards, retaining this amplification,
so that every surviving vertex at level \(j\)
has at least
\[
                   (m/m_j)^{2s}m^{-o_i(1)}
\]
terminal descendant time bins by low paths. In the pruned graph,
and after any subsequent masks, each root reaches at most
\begin{equation}\label{eq:source-39}
                           m_j^{2s}m^{o_i(1)}
\end{equation}
distinct time bins at level \(j\).
\end{lemma}

\begin{proof}
Select dyadic classes for all the indicated positive statistics.
For each time bin touched by a remaining parent \(v\), select
one remaining low child \(b\). Its intermediate neighborhood
contributes at least a constant times \(s_j'\) to \(D_j(v)\).
These contributions for distinct time bins are disjoint, because
each suffix has a unique child time bin. Thus the number of
touched bins is at most \(C d_j/s_j'\).

If \eqref{eq:source-38} failed by a fixed positive power of \(m\)
along a subsequence, the bound \(C d_j/s_j'\) would, after rebasing
the parent interval, be at most \(\ell_j^{2\sigma+o(1)}\) for some
fixed \(\sigma<s\). Here \(\log_m\ell_j\to1/J\), and the low radius
remains a fixed positive power in the \(\ell_j\) units within each
family. Thus the selected parent and child supports satisfy the
hypothesis of \eqref{ext:gap}. Its strict improvement for the full
transition contradicts the retained amplification, proving
\eqref{eq:source-38}.

For the pruning, begin at level \(J-1\) and proceed backwards.
At level \(j\), delete every parent touching fewer than
\(\ell_j^{2s}m^{-a}\) distinct child time bins in the current child
mask, initially with a fixed \(a>0\). The difference made to the
output is a diagram whose transition at this level starts from
the deleted parent set and ends in that current child mask.
That transition has the strict sparsity improvement just
described. A telescoping decomposition of the successive
mask changes contains only \(J\) such error diagrams. Each has
a fixed exponent gap, and the triangle inequality shows that
their total is negligible compared with a sufficiently accurate
near-extremal diagram.

We may consequently let the deficit \(a=a_i\downarrow0\) slowly
enough that its positive exponent gap still dominates the
amplification errors in the \(i\)-th family. This is a diagonal
choice: first use the gap for each fixed deficit, and then take
the initial accuracies and the family bases sufficiently far
out. No quantitative continuity of \(\gamma_\sigma\) at \(s\)
is required. The resulting pruned graph retains amplification
\(m^{\gamma-o_i(1)}\).

Each surviving parent now has at least
\(\ell_j^{2s}m^{-a_i}\) child time bins with surviving low children.
Choose one child in each such bin and iterate. The descendant
bins associated with distinct child time bins are disjoint.
Induction therefore gives at least
\[
 \prod_{h=j}^{J-1}\ell_h^{2s}m^{-a_i}
             =(m/m_j)^{2s}m^{-(J-j)a_i}
\]
terminal bins from every surviving level-\(j\) vertex.

Fix a root and one reachable vertex in each of its reachable
level-\(j\) time bins. The just constructed terminal descendants
are disjoint for these different time bins, and all are
accessible from the root in the original high graph. In that graph
this root reaches at most \(m^{2(s+\epsilon_i)}\) terminal time bins.
Dividing this bound by the lower descendant count gives
\eqref{eq:source-39}. If later masks remove some descendants,
the original pruned descendants remain valid witnesses in
this counting argument; later masks can only decrease the set
of reachable prefixes.
\end{proof}

\begin{lemma}[Incoming counts and terminal selection]\label{ext:leaves}
Further vertex masks can be chosen retaining amplification
\(m^{\gamma-o_i(1)}\) so that the following hold.
At each transition, every remaining child has between \(k_j\)
and \(2k_j\) low parents in the current parent mask, for a
positive integer \(k_j\). Put \(k=\prod_{j=0}^{J-1}k_j\).
Each terminal leaf has between \(k\) and \(2^Jk\) incoming low
paths, whereas any fixed root--leaf pair has at most
\(m^{o_i(1)}\) such paths. There is \(h_1>0\) such that
\[
                             |c_b|^2\asymp h_1k
\]
on all remaining leaves. Moreover, restricting the output to
any subset occupying a fraction \(m^{-o_i(1)}\) of these leaves
still retains amplification \(m^{\gamma-o_i(1)}\).
\end{lemma}

\begin{proof}
Proceed forwards. At transition \(j\), count the low parents of
each child in the already selected parent mask, discard the
zero-count class by the
\hyperref[ext:diagram-locality]{locality observation}, and select a positive
dyadic class retaining amplification \(m^{\gamma-o_i(1)}\). Subsequent forward
selections do not change any already fixed parent layer.
Thus the asserted incoming degrees hold in the final graph.
Induction from the roots gives the bounds \(k\) and \(2^Jk\)
for the number of incoming paths to each leaf. The bound for
a specified root--leaf pair follows from
Lemma~\ref{ext:endpoints}, because the radius powers tend to zero.

It remains to choose the leaves with the stronger norm property.
On the finite coordinate space, the lattice norm
\(\mathcal A\) has a nonnegative dual functional \(\lambda\) of
dual norm one satisfying
\[
                       \sum_b\lambda_b|c_b|=\mathcal A(c).
\]
One may take the absolute values of a norming functional:
the lattice property preserves its dual norm. Polynomial norm
comparisons and coordinate counts permit tiny coordinate amplitudes
to be discarded with negligible norm, hence negligible pairing with
\(\lambda\). Discarding additionally those contributions
\(\lambda_b|c_b|\) smaller than a sufficiently large negative power
of \(m\) times \(\mathcal A(c)\) loses negligible total pairing, by
the polynomial coordinate count. Partition the remaining leaves jointly into dyadic
classes for these two quantities. There are only a power of
\(\log m\) such classes. Some class \(\mathcal B\) satisfies
\[
 \sum_{b\in\mathcal B}\lambda_b|c_b|
       \geq (\log m)^{-C}\mathcal A(c),
 \qquad
 |c_b|^2\asymp h_1k\quad(b\in\mathcal B)
\]
for a suitable \(h_1>0\). Keep this class as a final leaf mask.
Incoming path counts for its leaves are unaffected.

If \(\mathcal B'\subset\mathcal B\) occupies a fraction \(f\),
comparability of \(\lambda_b|c_b|\) gives
\[
 \mathcal A(c|_{\mathcal B'})
 \geq \sum_{b\in\mathcal B'}\lambda_b|c_b|
 \geq \tfrac12 f\,(\log m)^{-C}\mathcal A(c).
\]
For \(f=m^{-o_i(1)}\), cubing this inequality loses only a
vanishing exponent. The functional need not norm any of the
later restricted outputs.
\end{proof}

\subsection{The conditional time law}

The masks now control prefix counts and incoming leaf degrees while
preserving near-extremal amplification. To apply geometric propagation,
we still need both support counts and maximal probabilities at every
time depth, conditional on each retained root. We obtain these from
uniform counting on a selected set of paths. This counting law is
auxiliary; the coefficients continue to come from the exact masked
matrix maps.

All graph statistics \(D_j\) and \(S_j\) still refer to the
original high graph. All low path counts in this subsection
refer to the final masks of Lemma~\ref{ext:leaves}.

\begin{lemma}[A linear conditional time profile]\label{ext:time-law}
There is a subset \(\Omega\) of the final low paths occupying
a fraction \(m^{-o_i(1)}\) of them such that uniform counting
on \(\Omega\) has these properties:
\begin{enumerate}[label=\textup{(\roman*)}]
\item Its root marginal is comparable to the uniform law on
a set of \(N_0\) roots.
\item Given each root, the terminal time satisfies
\eqref{eq:source-1}, with base \(M=m^2\), target exponent \(s\),
and error \(O(1/J)+o_i(1)\).
\end{enumerate}
The subset is used only to define this law; it does not change
any internal summand of the terminal coefficients.
\end{lemma}

\begin{proof}
Let \(H\) be the original high path count and \(L\) the final
low path count. By Lemma~\ref{ext:path-minimum},
\(L\geq m^{-\alpha_i}H\), where \(\alpha_i=o_i(1)\).
Write \(T^{(j)}\) for the level-\(j\) time interval containing
the terminal time.
For a vertex \(v\), let \(P_H(v)\) and \(P_L(v)\) be its high
and final low prefix counts, and let \(D_L(v)\) be its final
low suffix count. Thus \(P_L(v)\leq P_H(v)\) and
\(D_L(v)\leq D_j(v)\).

Choose \(\beta_i\downarrow0\) sufficiently slowly that
\(\beta_i-\alpha_i\) dominates the other errors and
\(m^{\alpha_i-\beta_i}\to0\) within each family. Call a remaining
level-\(j\) vertex bad if
\[
                         D_L(v)<m^{-\beta_i}D_j(v).
\]
The number of final low paths through bad vertices at this level
is bounded by
\[
 \sum_{v\ {\rm bad}}P_L(v)D_L(v)
 \leq m^{-\beta_i}\sum_v P_H(v)D_j(v)
 =m^{-\beta_i}H.
\]
In the final sum, \(v\) ranges over all vertices of the original
high graph at that level. Its equality to \(H\) holds because
every high path visits exactly one such vertex. Summing over
\(J+1\) levels shows that
the event \(\mathcal G\) that every visited vertex is good has
probability \(1-o(1)\) under uniform final low path counting.

Consider this unrestricted uniform low law, conditional on a
complete vertex prefix ending at a good \(v\) at level \(j\).
For a prescribed child time bin \(I\), its next-bin probability is
\begin{equation}\label{ext:one-bin}
 \frac{\displaystyle
       \sum_{\substack{b\ {\rm low\ child\ of}\ v\\I_b=I}}D_L(b)}
      {D_L(v)}
 \leq m^{\beta_i+o_i(1)}\frac{s_j'}{d_j}
 \leq m^{o_i(1)}\ell_j^{-2s}.
\end{equation}
If the numerator is nonzero, choose one remaining child \(b\)
in that bin. Its intermediate neighborhood contains every
other low child there, and the original high suffix sum on
the neighborhood is comparable to \(s_j'\). This bounds the
numerator. Goodness and \(D_j(v)\asymp d_j\) bound the
denominator, and \eqref{eq:source-38} gives the final inequality.

To iterate \eqref{ext:one-bin}, kill a path as soon as it
visits a bad vertex. Fix a nested sequence of time bins through
level \(j\). Conditional on each surviving \emph{vertex}
prefix, the next prescribed-bin probability is bounded by
\eqref{ext:one-bin}. Averaging over the vertex prefixes having
the same time prefix preserves that bound. Induction, with
the process killed at every step, therefore gives, for every
root \(v\) and level-\(j\) bin \(I\),
\begin{equation}\label{ext:killed}
 \Pr(\mathcal G,\ T^{(j)}=I\mid v)
       \leq m^{o_i(1)}\prod_{h<j}\ell_h^{-2s}
       =m^{o_i(1)}m_j^{-2s}.
\end{equation}
If the root is bad the left side is zero. In writing
\(\mathcal G\) on the left, one may first require survival only
through level \(j\), for which the induction applies directly,
and then impose survival at later levels. This can only
decrease the probability. Multiple vertex prefixes with the
same time prefix introduce no additional count.

Discard roots for which
\(\Pr(\mathcal G\mid v)<1/2\). Their total probability under
the original low root law is at most
\(2\Pr(\mathcal G^c)=o(1)\). Keep the good paths on the other
roots. Dividing \eqref{ext:killed} by their conditional survival
probability costs at most two, so the maximum probability of a
level-\(j\) time bin, conditional on each retained root, is
\(m^{-2s\log_m m_j+o_i(1)}\). The support bound at that level
is \eqref{eq:source-39}.

Between consecutive level depths, monotonicity of probabilities
uses the preceding level for the maximum weight, and monotonicity
of supports uses the following level for the number of occupied
bins. Since the gap in base-\(M=m^2\) depths is
\(1/J+o(1)\), these give \eqref{eq:source-1} at every dyadic
depth with error at most \(s/J+o_i(1)\).
The dyadic roundings change only bounded factors.

The retained measure is still uniform on a subset of paths.
Its positive root weights have polynomial ratios, because the
numbers of paths are positive integers with polynomial upper
bounds. Select a dyadic class of root weights carrying an
inverse-logarithmic fraction of the measure, and keep all
retained paths from those roots. If their number is \(N_0\),
the new root weights are comparable to \(1/N_0\).
This selection of whole root fibers does not change the
root-conditional time laws and retains a fraction
\(m^{-o_i(1)}\) of all final low paths.
\end{proof}

\subsection{Geometry bounds the capacity of the output}

Fix a positive \(\zeta<\gamma-10\kappa\).
Choose \(J\) sufficiently large that the \(O(1/J)\) error in
Lemma~\ref{ext:time-law} lies within the small-error range of
Theorem~\ref{geo:main} for exponent
\(1+2s-10\kappa-\zeta/2\), and then take \(i\) sufficiently large.
All the preceding selections are made for this fixed \(J\).

For a subset \(\mathcal P\) of the final low paths, say that a root
\(v\) is represented in a terminal bin \(I\) if some path of
\(\mathcal P\) starts at \(v\) and ends in \(I\). Count each such
root once in that bin, regardless of the number of its paths.
Call the paths in \(\mathcal P\) represented paths.

\Needspace{4\baselineskip}
\begin{lemma}[Represented roots]\label{ext:represented}
Fix any polynomial upper bound on the radii within each family.
There is a subset \(\mathcal P\), which may depend on this bound,
occupying a fraction \(m^{-o_i(1)}\) of all final low paths with the
following property. In every terminal bin \(I\), the number of
represented roots satisfying
\[
 U_{\nu_v}\in u+F,\qquad z_v+t_IU_{\nu_v}\in x+rF
\]
is at most
\begin{equation}\label{eq:source-41}
 K(d)\delta^{-2}W(F)\,
          m^{-2(1+s-10\kappa-\zeta)+o_i(1)},
 \qquad \text{radii}(F)\geq m^2\delta,
\end{equation}
uniformly over all independent centers \(u,x\in\R^2\) and all radii
up to the chosen upper bound.
One can keep a fraction \(m^{-o_i(1)}\) of the final leaves such
that every kept leaf has at least \(k m^{-o_i(1)}\)
represented incoming paths.
\end{lemma}

\begin{proof}
Put the uniform law of Lemma~\ref{ext:time-law} on its path
subset, take the root as base label, and assign it line parameters
\((U_{\nu_v},z_v)\). The root weights are comparable to \(1/N_0\),
and \(1\leq|d_v|\leq2\). The input capacity definition therefore
gives the density constant
\[
                         K_0\lesssim
                         \frac{K(d)\delta^{-2}}{N_0}
\]
for tests with radii at least \(\delta\).

We fix the width range before applying geometry. Choose a polynomial
cutoff \(H\geq m^2\delta\) bounding the norms of all tested velocities
and all spatial coordinates after division by \(r\), including
the root-line positions and terminal position labels. In the axes of
a rectangular test with half-widths \(e_j\), each coordinate interval
meets the bounded support in a set contained in an interval of
half-width \(\min(e_j,H)\). Choose the two centers independently,
using the same clipped half-widths and axes. The resulting joint test
covers the original test on the support, and its weight is no larger,
up to the fixed ellipse--rectangle comparison. Thus one polynomial
width range controls the full terminal capacity supremum.

Choose \(B=m^A\) larger than this cutoff, the upper width bound in
the statement, and the root coordinate bounds, with a fixed power
margin. This margin also covers the endpoint enlargements. With
\(M=m^2\), we have \(B\geq M\delta\). Rescale both line coordinates
by \(B^{-1}\).
The new input floor is
\[
 \delta'=\delta/B=M^{-D_m},\qquad
 D_m=\tfrac12(\log_mN+A)\geq1.
\]
The density constant becomes \(K_0B^2\), since \(W\) is
homogeneous of degree two. The output weight of a test is
\(W(F/B)=B^{-2}W(F)\), so these factors cancel on returning
to the old coordinates. The new output floor \(M\delta'\)
corresponds exactly to \(M\delta\) in the old coordinates.
This also covers \(m=N\), for which the unnormalized floor
would have given the inadmissible value \(D=1/2\).
To make the complexity fixed, take a subsequence with
\(D_m\to D_0\), and choose \(D_*=D_0+a\), where
\(0<a<\zeta/16\). Eventually
\(0<D_*-D_m<2a\). A radius at least \(M^{-D_*}\) but below
the available input floor \(M^{-D_m}\) can therefore be
rounded up by a factor at most \(M^{2a}\). Doing this to
both radii increases \(W\) by at most \(M^{4a}=m^{8a}\).
Thus the input cap at the fixed complexity \(D_*\) costs
at most this additional factor. The output floor for \(D_*\)
is smaller than the available output floor, so every required
output test is still allowed.

The definition of the geometric exponent as a small-error limit
of an infimum over finite complexities is used here: the
sufficiently small time-error threshold may be chosen for all
finite \(D_*\), although the large-base threshold and subpower
losses can depend on that fixed complexity. Thus the preceding
choice of \(J\), followed by \(i\) and then the family base,
is valid.

Apply Theorem~\ref{geo:main}, with base \(M=m^2\), to obtain a
dominated sublaw of subpower mass. The difference between the
chosen geometric exponent and \(1+2s-10\kappa-\zeta\) absorbs
the arbitrarily small fixed rounding loss just described.
After normalization, its
line-test probabilities satisfy
\begin{equation}\label{eq:source-40}
 \Pr\bigl(T_1=t,\ U\in F,\ X_t\in rF\bigr)
 \leq
 \frac{K(d)\delta^{-2}}{N_0}\,W(F)\,
 m^{-2(1+2s-10\kappa-\zeta)+o_i(1)} .
\end{equation}
As throughout, translations in the two tests are independent.

There are at most \(N_0m^{2s+o_i(1)}\) possible root--terminal-time
pairs. Discard pairs whose probability in this normalized law is
less than
\[
                           N_0^{-1}m^{-2s-\tau_i},
\]
where \(\tau_i=o_i(1)\) tends to zero slowly enough to dominate
the support error. This discards \(o(1)\) mass. Take \(\mathcal P\)
to be the path atoms in the support of the law after this deletion.
Each represented root in a specified terminal bin has pair probability
at least this threshold. Dividing \eqref{eq:source-40} by the threshold proves
\eqref{eq:source-41}, with all normalization factors absorbed
in \(m^{o_i(1)}\).

Let \(L\) again be the final low path count. The starting uniform
law for geometry is on a subset of at least
\(m^{-o_i(1)}L\) paths. Its dominated geometric sublaw retains
subpower mass. After normalization, the weight of any path
atom is consequently at most \(m^{o_i(1)}/L\).
The remaining support therefore represents at least
\(m^{-o_i(1)}L\) paths, even though their new weights need not
be equal.

Every original final leaf has between \(k\) and \(2^Jk\)
incoming low paths. Choose a threshold \(k m^{-\upsilon_i}\),
with \(\upsilon_i=o_i(1)\) sufficiently slow, and keep the
leaves with at least this many represented incoming paths.
The represented paths going to the other leaves are at most
\(k m^{-\upsilon_i}\) times the original leaf count, which is
negligible compared with the total represented path count.
Since no leaf has more than \(2^Jk\) incoming paths, the kept
leaves occupy a fraction \(m^{-o_i(1)}\) of all final leaves.
Lemma~\ref{ext:leaves} ensures that this last output restriction
retains amplification \(m^{\gamma-o_i(1)}\).
\end{proof}

\begin{lemma}[Terminal capacity]\label{ext:capacity}
For the kept coefficients of Lemma~\ref{ext:represented}, in
every terminal time bin one has
\begin{equation}\label{ext:capacity-bound}
 K(c_I|_{\mathrm{kept}})
 \leq h_1K(d)\,
          m^{2(1-s)+20\kappa+2\zeta+o_i(1)}.
\end{equation}
\end{lemma}

\begin{proof}
Consider an output test \(E,rE\), where the ordered radii of
\(E\) are \(e_1\geq e_2\geq\theta=m\delta\).
Let \(F\) have the same axes and radii
\[
                              f_j=\max(e_j,\delta/r).
\]
By Lemma~\ref{ext:endpoints}, the initial root of any represented
path ending in this output test belongs to the corresponding
line test \(F,rF\), enlarged by \(m^{o_i(1)}\).
Indeed the angular endpoint error is at most
\(m^{o_i(1)}\theta\), and the position error from the root line
is at most \(m^{o_i(1)}\delta\). The definition of \(F\) covers
both errors. Its radii satisfy the floor required in
\eqref{eq:source-41}. The harmless power enlargement changes
its weight and the resulting count only by \(m^{o_i(1)}\).

For a fixed represented root at this time, the terminal angular
grid supplies at most \(m^{o_i(1)}\) possible leaf angles.
The leaf positions lie both in the spatial test \(rE\) and in
an enlarged root beam of radius \(m^{o_i(1)}\delta\).
Lattice counting bounds their number by
\[
        m^{o_i(1)}w^{-2}\prod_{j=1}^2\min(re_j,\delta).
\]
No boundary correction is needed beyond a constant factor:
both \(re_j\) and \(\delta\) are at least the position mesh \(w\).
For any root and leaf, the number of paths is
\(m^{o_i(1)}\), by Lemma~\ref{ext:leaves}.

Each kept leaf has \(|c_b|^2\lesssim h_1k\) and at least
\(k m^{-o_i(1)}\) represented incoming paths. Its squared
coefficient is therefore bounded by \(h_1m^{o_i(1)}\) times
that incoming count. Summing represented paths in the test,
then using \eqref{eq:source-41}, gives
\begin{align*}
 K(c_I|_{\mathrm{kept}})
 &\leq m^{o_i(1)}
 \sup_E h_1\theta^2W(E)^{-1}
     w^{-2}\prod_{j=1}^2\min(re_j,\delta)\\
 &\hspace{25mm}\cdot
 K(d)\delta^{-2}W(F)
       m^{-2(1+s-10\kappa-\zeta)} .
\end{align*}
For the products of semiaxes,
\[
 \prod_{j=1}^2\min(re_j,\delta)
       =\delta^2\,\frac{e_1e_2}{f_1f_2},
 \qquad
 \frac{W(F)}{f_1f_2}
       =\left(\frac{f_1}{f_2}\right)^\kappa
       \leq\left(\frac{e_1}{e_2}\right)^\kappa
       =\frac{W(E)}{e_1e_2}.
\]
The inequality uses the reduction of aspect ratio when small
radii are raised to a common floor. Finally
\(\theta^2w^{-2}=r^{-2}=m^4\). Substitution proves
\eqref{ext:capacity-bound}.
\end{proof}

\subsection{Two mass bounds and the contradiction}

The square root of the capacity bound contributes
\(m^{1-s}h_1^{1/2}\), in addition to
\(m^{10\kappa+\zeta+o_i(1)}\). We bound the coefficient mass in
two ways: the frame estimate uses the number of times reached by each
root, and the fourth-power estimate gives a factor \(h_1^{-1}\).
Their minimum cancels \(m^{1-s}h_1^{1/2}\) for every value of \(h_1\).

\begin{lemma}[Terminal mass]\label{ext:mass}
The same kept output satisfies
\begin{equation}\label{ext:mass-bound}
 r w^2\sum_{b\ {\rm kept}}|c_b|^2
 \leq m^{o_i(1)}
       \min\bigl(m^{-2(1-s)},h_1^{-1}\bigr)
       \delta^2\#\mathcal R .
\end{equation}
\end{lemma}

\begin{proof}
For the first bound, fix a terminal time bin \(I\) and let
\(\mathcal R_I\) be the roots having an original high path to
that bin. Roots outside this set give a negligible error by the
\hyperref[ext:diagram-locality]{locality observation}, applied to the
original high graph with output restricted to \(I\).
Along the unique sequence of time ancestors of \(I\), the
successive exact masked transitions satisfy the squared-sum
bound of Lemma~\ref{pkt:frame}, with the parent-bin rebasing
of Lemma~\ref{pkt:composition}. The transition at level \(j\)
has squared operator norm at most \(\ell_j^2\), and masks are
orthogonal coordinate projections. With the scale
normalizations this gives
\[
                    w^2\sum_{b\in I}|c_b|^2
                    \lesssim \delta^2\#\mathcal R_I
\]
apart from negligible error. Summation with weight \(r=m^{-2}\)
counts each root at most \(m^{2(s+\epsilon_i)}\) times by the
original high time-support bound. This proves the first term in
\eqref{ext:mass-bound}.

For the second bound, apply Proposition~\ref{pkt:quartic} at each
actual masked transition. Lemma~\ref{ext:leaves} bounds its
low-parent multiplicity by \(2k_j\). Use
Remark~\ref{pkt:vanishing-powers} with adjacency power
\(\eta_i=\eta_i^{\mathrm{lo}}\): within each family the localization
and sampling powers are chosen below this fixed positive power.
Thus the full masked transition has the required multiplicity bound,
and the product of its localization, coloring, and decoupling losses
over the fixed \(J\) transitions is \(m^{o_i(1)}\).
The factor \(2^J\) is absorbed in the same loss.
Successive applications of \eqref{eq:source-33} yield
\[
 r w^2\sum_{b\ {\rm kept}}|c_b|^4
 \leq k\,m^{o_i(1)}
          \delta^2\sum_{v\in\mathcal R}|d_v|^4
 \lesssim k\,m^{o_i(1)}\delta^2\#\mathcal R .
\]
Since \(|c_b|^2\asymp h_1k\) on the kept leaves, division by
\(h_1k\) proves the second term in \eqref{ext:mass-bound}.
\end{proof}

\begin{proof}[Completion of the proof of Theorem~\ref{pkt:main}]
For each terminal bin, use its entire kept coefficient array
as a single atom in the definition of \(A\). Equations
\eqref{ext:capacity-bound} and \eqref{ext:mass-bound} imply
\begin{align*}
 \mathcal A(c|_{\mathrm{kept}})^3
 &\leq r\sum_I G(c_I|_{\mathrm{kept}})^3\\
 &\leq m^{10\kappa+\zeta+o_i(1)}
     \delta^2\#\mathcal R\,K(d)^{1/2}\\
 &\hspace{10mm}\cdot
   m^{1-s}h_1^{1/2}
      \min\bigl(m^{-2(1-s)},h_1^{-1}\bigr).
\end{align*}
For every positive \(h_1\),
\[
 h_1^{1/2}\min\bigl(m^{-2(1-s)},h_1^{-1}\bigr)
                      \leq m^{-(1-s)}.
\]
Also \(1\leq|d_v|\leq2\) gives
\(G(d)^3\asymp\delta^2\#\mathcal R\,K(d)^{1/2}\).
The cubed amplification of the kept output is consequently
at most \(m^{10\kappa+\zeta+o_i(1)}\).

By Lemmas~\ref{ext:leaves} and \ref{ext:represented}, this same
output has amplification at least \(m^{\gamma-o_i(1)}\).
Our choice \(\zeta<\gamma-10\kappa\), followed by sufficiently
large \(i\) and then sufficiently large \(m\) within that
family, is a contradiction. Thus the assumption
\(\gamma>10\kappa\) is impossible, which proves the propagation
bound and its uniform finite-family formulation.
\end{proof}
